\section{Discussion and Analysis}
\label{sec:discussion}

\subsection{Edge Feasibility and Hardware Profiling}
A central design goal of Dioptra-DINO is edge feasibility. Most spatial AI models assume cloud GPU servers or heavy onboard workstations (NVIDIA RTX 4090/A100). As documented in Section~\ref{sec:experiments}, Dioptra-DINO executes steady-state forward inference in \textbf{$25.92$\,ms} ($\approx \mathbf{38.6}$\,FPS) on consumer Apple Silicon (Metal Performance Shaders, unquantized FP32) and in \textbf{$21.8$\,ms} ($\approx \mathbf{45.9}$\,FPS) on an NVIDIA Tesla T4 GPU (PyTorch 2.1, CUDA 12.1, FP32 unquantized, batch size 1). The 14-pixel patch tokenization of DINOv2-Small ($16 \times 16 = 256$ tokens) drastically suppresses quadratic self-attention memory, operating with $9.4\times$ less attention compute than patch-8 baselines.

Crucially, Dioptra-DINO's single-image inference consumes under \textbf{$210$\,MB of RAM}, while its complete model parameter footprint is only \textbf{$105.05$\,MB} in FP32 format ($52.5$\,MB in FP16). The lightweight parameter profile enables seamless deployment on micro-aerial vehicles (MAVs), field rovers, and mobile edge robotics.

\subsection{Qualitative Surface Reconstruction \& Normal Consistency}
As shown in Figure~\ref{fig:teaser}, depth predictions from Dioptra-DINO produce consistent 3D surface geometry. Surface normals derived via spatial cross products reveal that the network distinguishes horizontal upward-facing ground planes from orthogonal vertical building walls and pillars. Multi-Scale Virtual Normal Loss (VNL) supervision achieves a median 3D surface normal error of $39.80^\circ$ ($40.17^\circ$ mean), which is competitive for a compact $27.5$\,M parameter model without external normal estimation backbones.

\subsection{Failure Mode Analysis: Downward Close-ups}
Analyzing challenging frames (e.g., frames with $\text{AbsRel} > 1.5$) reveals distinct failure modes common to monocular metric depth:
\begin{enumerate}
    \item \textbf{Foreground Dominance}: When the camera tilts steeply downward facing an obstacle or catwalk within $<0.8$\,m that occupies $>60\%$ of the visual field, the image median anchors strictly to the near plane.
    \item \textbf{Median Normalization Skew}: Under standard evaluation protocols, dividing by an extreme near-field median mathematically amplifies relative error on the small sliver of distant background through dividing by true distance.
    \item \textbf{Specular and Textureless Highlights}: Extreme specular reflection on flat metallic surfaces eliminates vanishing lines and textural gradients, requiring temporal stereo or visual-inertial odometry fusion for disambiguation.
\end{enumerate}

\subsection{Camera-Intrinsic Equivariance vs. Positional Embedding}
\label{sec:fov_discussion}
A critical theoretical question for ray-conditioned networks~\cite{facil2019camconvs,yin2023metric3d,piccinelli2024unidepth} is whether explicit camera parameters $\mathbf{K}$ provide actionable geometric guidance or merely act as a passive coordinate bias. As demonstrated in Section~\ref{sec:fov_and_stratification}, when trained exclusively on fixed-sensor benchmarks such as TartanAir, models can overfit to the static relationship between image coordinates $(u, v)$ and ray vectors $\hat{\mathbf{r}}(u, v)$, rendering the ray embedding functionally equivalent to standard sinusoidal positional encoding. 

Our dynamic pinhole augmentation uncouples pixel space from metric ray directions by introducing stochastic window cropping ($L \in [0.55\cdot S_{\max}, S_{\max}]$) and continuous focal/center recalculations. As demonstrated by the error reductions across multi-camera field-of-view sweeps (Table~\ref{tab:multi_fov_and_stratification}), conditioning depth decoding on dynamically varying $\hat{\mathbf{r}}(u, v)$ encourages the network to learn optical equivariance:
\begin{equation}
    f_\theta(\mathbf{I} \circ \mathcal{T}_s, \mathbf{K}') \approx f_\theta(\mathbf{I}, \mathbf{K})
\end{equation}
where physical 3D scene distances remain invariant under optical focal scaling. The dynamic pinhole formulation substantially mitigates scale drift across varying optical configurations (improving the scale calibration ratio to $1.192$ vs. $1.593$ for the canonical 2D ViT). While wide-angle unprojection yields pronounced error drops (up to $-69.1\%$) on structured architectural scenes, performance slightly degrades ($+8.2\%$) in foliage-dense unstructured scenes, illustrating that geometric ray unprojection is most effective when scenes contain clear vanishing cues.

\subsection{Limitations and Experimental Scope}
While Dioptra-DINO demonstrates the viability of camera-grounded geometric embeddings in a compact model ($27.5$\,M parameters), several limitations delineate its experimental scope:
\begin{itemize}
    \item \textbf{Architectural Context vs. Head-to-Head Comparison}: While Table~\ref{tab:env_breakdown} includes foundation models (ZoeDepth, Metric3D) to provide orders-of-magnitude context on model size, compute, and reported literature metrics, these models were trained on multi-dataset collections. Our primary empirical claims are strictly grounded in our controlled 7-model ablation matrix within TartanAir, where all variants are trained and evaluated on identical splits with identical protocols.
    \item \textbf{Controlled Ablation Progression}: As demonstrated in our progression analysis (Section~\ref{sec:ablations} and Table~\ref{tab:ablation_progression}), isolating the geometric attention bias, ray positional encoding, and camera conditioning confirms that ARA improves normal planarity by $3.21^\circ$, while full DINOv2 integration doubles raw $\delta_1$ accuracy from $17.4\%$ to $35.8\%$ in unaligned metric space.
    \item \textbf{Cross-Domain Real-World Transfer}: The present evaluation is conducted on photorealistic TartanAir simulation trajectories. Evaluating zero-shot transfer onto physical real-world benchmarks (such as NYUv2, KITTI, and real robot platforms) with multi-thousand-frame held-out distributions represents an important direction for physical robotics deployment.
    \item \textbf{Quantized Edge Acceleration}: FP32 inference achieves $45.9$\,FPS on NVIDIA Tesla T4 and $38.6$\,FPS on Apple Silicon GPU. Porting model weights to FP16 CoreML and INT8 TensorRT backends represents a direct path to exceeding $>60$\,FPS on embedded robotics platforms.
\end{itemize}
